\documentclass[11pt,reqno]{amsart}
\usepackage[margin=1in]{geometry}
\usepackage{amsmath,amssymb,amsthm,mathtools}
\usepackage{enumitem}
\usepackage{microtype}
\usepackage{hyperref}

\theoremstyle{plain}
\newtheorem{theorem}{Theorem}
\newtheorem{lemma}[theorem]{Lemma}
\newtheorem{proposition}[theorem]{Proposition}
\newtheorem{corollary}[theorem]{Corollary}
\newtheorem{conjecture}[theorem]{Conjecture}

\theoremstyle{definition}
\newtheorem{definition}[theorem]{Definition}
\newtheorem{remark}[theorem]{Remark}
\newtheorem{example}[theorem]{Example}

\newcommand{\Z}{\mathbb{Z}}
\newcommand{\Q}{\mathbb{Q}}
\newcommand{\F}{\mathbb{F}}
\newcommand{\Zn}{\Z/n\Z}
\newcommand{\Zns}{(\Z/n\Z)^{\times}}
\newcommand{\Zts}{\Z/10\Z}
\newcommand{\meet}{\wedge}
\newcommand{\Tstar}{T^{*}}
\newcommand{\disc}{\mathrm{disc}}
\newcommand{\DYN}{\mathrm{DYN}}
\newcommand{\sinc}{\mathrm{sinc}}
\newcommand{\ord}{\mathrm{ord}}
\newcommand{\Maj}{\mathrm{Maj}}

\title[The forced $5/7$ torus (up to a calibration choice)]{The Forced $5/7$ Torus Aspect Ratio (Up to a Calibration Choice):\\
Cyclotomic Forcing on $\Z/10\Z$}

\author{Brayden R.\ Sanders}
\address{7Site LLC, Hot Springs, Arkansas, USA}
\email{brayden@7site.co}

\author{M.\ Gish}
\address{Independent Researcher}

\date{\today}

\begin{document}

\begin{abstract}
The ring $\Zts$ carries four simultaneous algebraic structures --- additive divisor chain, multiplicative orbit lattice, additive translation flow $x \mapsto x+1$, and multiplicative root flow $x \mapsto gx$. The additive structure is totally ordered by divisor refinement; the multiplicative structure has pairwise incompatible CRT-factor partitions for distinct primes (Sanders--Mayes \emph{UOP}, Lemma~4.1), so the two are not jointly orderable. The companion Sanders--Gish \emph{Flatness Theorem} (submitted to \emph{J.\ Pure Appl.\ Algebra}) shows that the minimal smooth $2$-manifold simultaneously hosting all four structures is the torus $T^2 = S^1 \times S^1$. We isolate the cyclotomic data that determines the aspect ratio of this torus and prove the following structural statement: \emph{relative to the cyclotomic embedding scale, in which a prime-$p$ closed circle has circumference $p$,} the aspect ratio $T^{*} := R/r$ equals $5/7$, where $R$ is the smallest prime $p \mid n$ at which the cyclotomic value $A_p = 2 \cos(\pi/p)$ has algebraic degree $\le 2$ over $\Q$ (giving $p = 5$, with $A_5 = \varphi$, the golden ratio) and $r$ is the smallest prime $p$ at which $A_p$ has algebraic degree $\ge 3$ over $\Q$ (giving $p = 7$, with minimal polynomial $x^{3} - x^{2} - 2x + 1$, irreducible of degree $3$). The forcing is conditional on the cyclotomic-embedding calibration; we are explicit about this dependence and discuss what would be required to remove it (\S\ref{sec:torus} Remark~\ref{rem:calibration}). The numerical value $5/7$ arises in several further ways inside the authors' research program; we sketch them in \S\ref{sec:companions}, distinguishing the genuinely independent derivations from those that are reformulations of the present cyclotomic derivation.
\end{abstract}

\maketitle

\section*{Lens and substrate}

This paper works on the substrate $\Zts = \Z/10\Z$ with the four ring structures of Definition~\ref{def:four-structures}. The choice of $\Zts$ (as opposed to $\Z/N$ for other squarefree $N$, or a finite field $\F_p$) reflects a structural reading of the substrate motivated by the four-flow decomposition of the companion \emph{Flatness Theorem} (Sanders--Gish), not by a first-principles derivation. The theorems below are theorems on this specific substrate; the generalization conjecture in \S\ref{sec:gen} states what we expect on the analogous substrates $\Z/n$ for squarefree $n$ with $5 \mid n$ and $n > 5$. We are also explicit that Theorem~\ref{thm:forced57} is conditional on the cyclotomic-embedding calibration of Definition~\ref{def:cyclo-cal}; this calibration is itself imported from the \emph{Flatness Theorem} rather than derived here (Remark~\ref{rem:calibration}). The substrate-and-calibration scope of the result is exactly the structural reading of the cyclotomic threshold described in the abstract; whether other substrates or other calibrations produce a similarly rigid threshold is open. Domain precedent for small finite commutative non-associative structures is Dr\'apal and Wanless \cite{DrapalWanless} (opposite extremum: maximally non-associative).

The verification script \texttt{verify\_J13.py} (accompanying this submission) reproduces, at machine precision, the minimal-polynomial identification, the rational-root irreducibility test, the discriminant $\disc(g) = 49 = 7^{2}$, the Galois group $A_{3} \cong \Z/3\Z$ (by the disc-square criterion for cubics), and the cyclotomic degree threshold $\deg_{\Q}(A_{p}) = (p-1)/2$ at the first four odd-prime values $p = 2, 3, 5, 7$. The script is sympy-based, runs in under five seconds, and PASSes 6/6.

\section{Introduction}

\subsection{Setup}

Let $n \geq 6$ be squarefree with $k \geq 2$ distinct prime factors, $n = p_1 \cdots p_k$. The Chinese Remainder Theorem provides $\Zn \cong \prod_i \Z/p_i\Z$ and $\Zns \cong \prod_i (\Z/p_i\Z)^{\times}$. We focus on the smallest non-trivial squarefree case, $n = 10 = 2 \cdot 5$, where the forced torus aspect ratio takes its canonical value $5/7$.

The cyclotomic numbers
\[
A_p := 2 \cos(\pi/p) \qquad (p \text{ prime}, \, p \geq 2)
\]
arise as the trace of the matrix $T_g$ (multiplication by a primitive root $g$ modulo $p$) acting on the cyclotomic embedding $\Phi_p \colon \Z/p\Z \to S^1$, $x \mapsto e^{2\pi i x/p}$. Their algebraic structure over $\Q$ controls the algebraic structure of the multiplicative flow's orbit decomposition.

\begin{remark}
The algebraic degrees of the first cyclotomic values over $\Q$ are
\[
A_2 = 0,\quad A_3 = 1,\quad A_5 = \varphi = (1+\sqrt 5)/2 \in \Q(\sqrt 5),\quad A_7 \in \Q[x]/(x^3 - x^2 - 2x + 1).
\]
The minimal polynomial of $A_p = 2\cos(\pi/p)$ over $\Q$ has degree $\varphi(2p)/2 = (p-1)/2$ for $p$ an odd prime (Lehmer, 1933; Watkins--Zeitlin, 1993), so for $p = 2,3,5,7$ the degrees are $0, 1, 2, 3$. The first prime where $A_p$ has degree $\ge 3$ over $\Q$ is therefore $p = 7$. (Caution: the polynomial $8x^{3} - 4x^{2} - 4x + 1$, which appears in many references, is the minimal polynomial of $\cos(\pi/7)$, not of $A_7 = 2\cos(\pi/7)$; the two are related by the substitution $x \mapsto x/2$.)
\end{remark}

\subsection{Companion to the Flatness Theorem}

The Sanders--Gish \emph{Flatness Theorem} (companion, submitted to \emph{J.\ Pure Appl.\ Algebra}) shows that the four ring structures on $\Z/10\Z$ cannot be simultaneously embedded in a flat $2$-dimensional surface, and that the minimal smooth surface admitting the embedding is a torus $T^2 = S^1 \times S^1$. The present paper takes the existence of the torus as input from the Flatness Theorem and determines its aspect ratio.

\subsection{Main result}

\begin{theorem}[Cyclotomic-calibrated $5/7$ aspect ratio] \label{thm:forced57}
Fix the cyclotomic-embedding calibration of Definition~\ref{def:cyclo-cal} below: the major circle $S^{1}_{R}$ inherits the closed-cycle circumference of the additive flow on the prime-$p$ component of the CRT decomposition with $p$ the smallest prime divisor of $n$ at which $A_{p}$ is irrational and $\deg_{\Q}(A_{p}) \le 2$ (excluding the degenerate cases $A_{p} \in \{0, 1\}$, i.e.\ $p = 2$ and $p = 3$). The minor circle $S^{1}_{r}$ inherits the analogous circumference at the smallest prime $p$ (not necessarily dividing $n$) at which $\deg_{\Q}(A_{p}) \ge 3$. Under this calibration, for $n = 10$,
\[
\Tstar := R/r = 5/7.
\]
\end{theorem}

The proof is structural: $R \propto 5$ arises because $5$ is the smallest prime divisor of $10$ at which the cyclotomic value $A_p$ admits a quadratic algebraic relation (giving the identity $A_5 = \varphi \in \Q(\sqrt{5})$), while $r \propto 7$ arises because $7$ is the smallest prime at which $A_p$ first acquires an irreducible cubic minimal polynomial. Both proportionality assignments are conditional on the cyclotomic-embedding calibration; \emph{the calibration is itself a structural choice imported from the Flatness Theorem and is not derived from the present paper.} See Remark~\ref{rem:calibration} for what a derivation of the calibration would require.

\section{The four structures and the torus} \label{sec:torus}

We summarize the structural inputs from the companion Flatness Theorem.

\begin{definition}[The four structures on $\Zn$] \label{def:four-structures}
\begin{enumerate}[label=(\arabic*)]
\item \emph{Additive structure} ($A$-Struct): the divisor chain of residue partitions $\{\pi_d : d \mid n\}$. Totally ordered by refinement; isomorphic to the divisor lattice of $n$ (flat).
\item \emph{Multiplicative structure} ($M$-Struct): the orbit lattice of partitions $\{\pi_{\DYN}(G) : G \leq \Zns\}$. For squarefree $n$ with $k \geq 2$, the CRT factor partitions $\pi_{p_i}, \pi_{p_j}$ are pairwise incompatible (Sanders--Mayes \emph{UOP}, Lemma~\ref{lem:pwise-incomp-cite}). Not totally ordered.
\item \emph{Additive flow} ($A$-flow): the dynamics of $x \mapsto x + 1$ on $\Zn$, generating one closed cycle of length $n$.
\item \emph{Multiplicative flow} ($M$-flow): the dynamics of $x \mapsto g x$ for a generator $g \in \Zns$. Orbits have lengths dividing $\ord_n(g) \mid \varphi(n)$.
\end{enumerate}
\end{definition}

\begin{lemma} \label{lem:pwise-incomp-cite}
For squarefree $n$ and distinct primes $p_i, p_j \mid n$, $\pi_{p_i}$ and $\pi_{p_j}$ are incompatible. \emph{(Sanders--Mayes, UOP companion, Lemma 4.1.)}
\end{lemma}

\begin{theorem}[Flatness obstruction]\label{thm:flatness-cite}
The four structures cannot be simultaneously embedded in a flat $2$-dimensional surface; the minimal smooth surface admitting the joint embedding is $T^2 = S^1 \times S^1$. \emph{(Sanders--Gish, Flatness Theorem, Theorems 1 and 2.)}
\end{theorem}

The two $S^1$ factors of $T^2$ correspond respectively to:
\begin{itemize}
\item \emph{Major circle} ($S^1_R$): hosts the $A$-flow. Circumference $\propto R$, where $R$ is determined by additive cyclotomic closure (\S\ref{sec:R-major}).
\item \emph{Minor circle} ($S^1_r$): hosts the $M$-flow. Circumference $\propto r$, where $r$ is determined by multiplicative cyclotomic obstruction (\S\ref{sec:r-minor}).
\end{itemize}

\begin{definition}[Cyclotomic-embedding calibration] \label{def:cyclo-cal}
Fix the embedding $\Phi \colon \Z/n\Z \to \prod_{i} S^{1}$ given by $\Phi(x) = (e^{2\pi i x_{i} / p_{i}})_{i}$ where $x_{i}$ is the image of $x$ in the CRT factor $\Z/p_{i}\Z$. We call this the \emph{cyclotomic-embedding calibration}. In this calibration, the closed-cycle on the prime-$p$ component has circumference exactly $p$ (one revolution per CRT-factor cycle). The torus aspect ratio depends on this calibration: a different choice of embedding rescales $R$ and $r$ independently and would change their ratio. The Flatness Theorem (Sanders--Gish, companion) selects this calibration as the one in which the four ring structures embed isometrically; a self-contained derivation of the canonical calibration directly from $\Z/n\Z$ would yield an unconditional version of Theorem~\ref{thm:forced57} and is left to future work.
\end{definition}

\begin{remark}[On the conditional nature of Theorem~\ref{thm:forced57}] \label{rem:calibration}
Theorem~\ref{thm:forced57} establishes that, \emph{within the cyclotomic-embedding calibration}, the aspect ratio is the cyclotomic threshold ratio $5/7$. We do not claim that $5/7$ is independent of the calibration: a torus admits a one-parameter family of embeddings differing in the relative scale of the two circle factors, and only the cyclotomic-embedding calibration of Definition~\ref{def:cyclo-cal} (in which a prime-$p$ closed circle has circumference $p$) yields the specific value $5/7$. The Flatness Theorem singles out this calibration on structural grounds (it is the unique calibration in which the additive and multiplicative flows close into circles of integer circumference), but our paper takes that selection as input rather than re-deriving it. We therefore use ``forced'' throughout in the qualified sense ``forced relative to the cyclotomic-embedding calibration.''
\end{remark}

\section{The major radius: additive cyclotomic closure} \label{sec:R-major}

\begin{theorem}[Major-radius selection under cyclotomic calibration] \label{thm:R-forced}
Under the cyclotomic-embedding calibration of Definition~\ref{def:cyclo-cal}, the major-circle circumference $R$ equals the smallest prime $p \mid n$ at which the cyclotomic value $A_p = 2\cos(\pi/p)$ has algebraic degree $\leq 2$ over $\Q$ and exceeds the degenerate cases $A_p \in \{0, 1\}$. For $n = 10$, $R = 5$.
\end{theorem}

\begin{proof}
The $A$-flow is the cyclic translation $x \mapsto x+1$ on $\Zn$, with single orbit of length $n$. By CRT, $\Zn \cong \prod_i \Z/p_i\Z$, and the embedding $\Phi$ factors through prime-component embeddings $\Phi_{p_i}(x_i) = e^{2 \pi i x_i / p_i}$. Multiplication by a primitive $p_i$-th root of unity $\zeta_{p_i}$ acts on $\Phi_{p_i}(\Z/p_i\Z)$ as the additive shift, with trace
\[
\mathrm{Tr}(M_{\zeta_{p_i}}) = \zeta_{p_i} + \zeta_{p_i}^{-1} = 2 \cos(2\pi/p_i).
\]
The corresponding cyclotomic value at the half-angle is $A_{p_i} = 2 \cos(\pi/p_i)$, and the trace structure of the multiplicative action on $\Phi_{p_i}$ is governed by $\Q(A_{p_i})$.

\emph{Closure condition.} Within the cyclotomic-embedding calibration, the $A$-flow on the $p_i$-component encodes a $p_i$-circumference closed circle whose embedding extends from $\Q$ through the quadratic field $\Q(A_{p_i})$ exactly when $\deg_{\Q}(A_{p_i}) \le 2$. When $A_{p_i} \in \Q$ (degenerate case: $A_2 = 0$, $A_3 = 1$), the prime contributes only a degree-$0$ or degree-$1$ closure, and the resulting circle is structurally collapsed (we discuss this in Remark~\ref{rem:degenerate} below). The first non-degenerate quadratic closure occurs at the smallest prime $p \mid n$ with $A_p$ irrational and $\deg_{\Q}(A_p) = 2$.

\emph{Compute.} $A_2 = 0$, $A_3 = 1$ (degenerate); $A_5 = (1+\sqrt{5})/2 = \varphi \in \Q(\sqrt{5})$, degree $2$ (first non-degenerate). Among the prime factors of $n = 10 = 2 \cdot 5$, the smallest with non-degenerate quadratic closure is $p = 5$, so under the cyclotomic-embedding calibration, $R = 5$.
\end{proof}

\begin{remark}[Degenerate primes] \label{rem:degenerate}
At $p = 2$, the cyclotomic value $A_2 = 0$ is the trivial root of $x = 0$; the prime-$2$ component of the CRT decomposition embeds into the two-point set $\{1, -1\} \subset S^{1}$, which is a degree-$1$ object. At $p = 3$, the value $A_3 = 1$ is rational; the prime-$3$ component embeds into a non-degenerate triangle in $S^{1}$ but the cyclotomic field $\Q(A_3)$ is just $\Q$. We exclude these from the major-radius selection because the $A$-flow's closed circle requires a non-trivial cyclotomic structure --- the very algebraic relation $A_5 = \varphi$ that distinguishes a closed pentagon from a discrete five-point set. This exclusion is structural and is the reason the conjecture in \S\ref{sec:gen} requires care for small $n$.
\end{remark}


\section{The minor radius: multiplicative cyclotomic obstruction} \label{sec:r-minor}

\begin{theorem}[Minor-radius selection under cyclotomic calibration] \label{thm:r-forced}
Under the cyclotomic-embedding calibration of Definition~\ref{def:cyclo-cal}, the minor-circle circumference $r$ equals the smallest prime $p$ at which the cyclotomic value $A_p$ has algebraic degree $\geq 3$ over $\Q$, equivalently the smallest prime where the minimal polynomial of $A_p$ over $\Q$ is irreducible of degree $3$ or higher. This prime is $p = 7$, so $r = 7$.
\end{theorem}

\begin{proof}
The $M$-flow is multiplication by a primitive root $g \in \Zns$, generating orbits of lengths dividing $\ord_n(g) \mid \varphi(n)$. The continuum limit of the $M$-flow on the cyclotomic embedding produces standing harmonic waves --- the $\sinc^2$ resonance field of \cite{SandersGish-FirstG}, with values
\[
\mathcal{R}(k, f) = \frac{\sin^2(\pi k f)}{k^2 \sin^2(\pi f)}.
\]
For $f = 1/p$, $\mathcal{R}(k, 1/p) = 0$ iff $p \mid k$ (\cite{SandersGish-FirstG}, Theorem~1). The harmonic structure of the $M$-flow is therefore controlled, prime by prime, by the cyclotomic numbers $A_p = 2\cos(\pi/p)$.

\emph{Obstruction condition.} Within the cyclotomic-embedding calibration, the $M$-flow's harmonic interference pattern on the $p$-component factors through the cyclotomic field $\Q(A_p)$ as a (at most) quadratic extension precisely when $\deg_{\Q}(A_p) \le 2$; when $\deg_{\Q}(A_p) \ge 3$, the $M$-flow's resonance structure cannot be expressed by a quadratic algebraic identity in the cyclotomic embedding, and we say the prime $p$ \emph{obstructs} the closure.

\emph{Compute.} The minimal polynomial of $A_p = 2\cos(\pi/p)$ over $\Q$ has degree $\varphi(2p)/2 = (p-1)/2$ for $p$ an odd prime (Lehmer 1933; Watkins--Zeitlin 1993). Concretely:
\begin{itemize}
\item $p = 2$: $A_2 = 0$, degree $0$, no obstruction (degenerate).
\item $p = 3$: $A_3 = 1$, minimal polynomial $x - 1$, degree $1$, no obstruction (degenerate).
\item $p = 5$: $A_5 = \varphi$, minimal polynomial $x^2 - x - 1$, degree $2$ (closure within $\Q(\sqrt 5)$).
\item $p = 7$: $A_7 = 2 \cos(\pi/7)$, minimal polynomial $x^3 - x^2 - 2x + 1$ over $\Q$, irreducible of degree $3$ (Lemma~\ref{lem:A7-irreducible} below). This is the first cubic obstruction.
\end{itemize}

So the smallest prime at which the multiplicative cyclotomic flow first encounters a cubic-or-higher obstruction is $p = 7$, and under the cyclotomic-embedding calibration $r = 7$.

\emph{Note on $7 \nmid 10$.} The obstruction prime is determined by the global cyclotomic structure of $\Q$ rather than by the prime divisors of $n$ alone. The harmonic resonance field $\mathcal{R}(k, 1/p)$ of the $M$-flow makes sense for every prime $p$ (as a function on $\Z$), and the smallest $p$ at which $\Q(A_p)$ is non-quadratic is $p = 7$ regardless of whether $7 \mid n$. Within the cyclotomic-embedding calibration, this universal obstruction is what fixes the minor-circle circumference.
\end{proof}

\begin{lemma}[Irreducibility of the $A_7$ minimal polynomial] \label{lem:A7-irreducible}
The polynomial $g(x) = x^{3} - x^{2} - 2x + 1$ is the minimal polynomial of $A_7 = 2\cos(\pi/7)$ over $\Q$, and is irreducible over $\Q$.
\end{lemma}

\begin{proof}
\emph{Identification of the polynomial.} The Chebyshev identity $U_{p-1}(\cos\theta) = \sin(p\theta) / \sin\theta$ at $\theta = \pi/p$ gives $\sin(p \cdot \pi/p) = \sin\pi = 0$, so $\cos(\pi/p)$ is a root of $U_{p-1}(x)$, and consequently $2\cos(\pi/p)$ is a root of the rescaled polynomial $\widetilde{U}_{p-1}(x) := 2^{p-1} U_{p-1}(x/2) / x$ (when $p$ is odd, the relevant factor of $\widetilde{U}_{p-1}$ is the degree-$(p-1)/2$ minimal polynomial). For $p = 7$, the standard reference (Lehmer 1933; Watkins--Zeitlin 1993, Theorem~3) gives the minimal polynomial of $2\cos(\pi/7)$ over $\Q$ as $g(x) = x^{3} - x^{2} - 2x + 1$. (Caution: the polynomial $8x^{3} - 4x^{2} - 4x + 1$ is the minimal polynomial of $\cos(\pi/7)$, not of $2\cos(\pi/7)$; substituting $x \mapsto x/2$ in the former and clearing denominators yields the latter.) Direct check: the three roots of $g$ are $\{2\cos(\pi/7), 2\cos(3\pi/7), 2\cos(5\pi/7)\}$, which are conjugate under the Galois action of $\mathrm{Gal}(\Q(\zeta_{14})/\Q)$.

\emph{Irreducibility.} A monic cubic over $\Z$ is reducible over $\Q$ iff it has a rational root, and by the rational root theorem any rational root of a monic integer polynomial is an integer dividing the constant term. The only candidates are $\pm 1$. Compute:
\[
g(1) = 1 - 1 - 2 + 1 = -1 \neq 0, \qquad g(-1) = -1 - 1 + 2 + 1 = 1 \neq 0.
\]
Hence $g$ has no rational roots and is irreducible over $\Q$.
\end{proof}

\begin{remark}[Galois group and discriminant]
The discriminant of $g(x) = x^{3} - x^{2} - 2x + 1$ is $\disc(g) = 49 = 7^{2}$, a perfect square. Therefore $\mathrm{Gal}(g/\Q) \cong A_{3} \cong \Z/3\Z$ (the cyclic group of order $3$, not the full symmetric group $S_3$), consistent with the classical fact that the splitting field of $g$ is the totally real cubic subfield of $\Q(\zeta_{7})$. In particular, the obstruction is structural: $\Q(A_{7})$ is a Galois cubic extension of $\Q$, not merely a non-normal cubic.
\end{remark}

\section{Aspect ratio: $R/r = 5/7$}

\begin{proof}[Proof of Theorem~\ref{thm:forced57}]
Apply the cyclotomic-embedding calibration of Definition~\ref{def:cyclo-cal}. Theorem~\ref{thm:R-forced} gives $R = 5$ (smallest prime $p \mid 10$ with non-degenerate $\deg_{\Q}(A_{p}) \le 2$), and Theorem~\ref{thm:r-forced} gives $r = 7$ (smallest prime, $p \mid n$ or not, with $\deg_{\Q}(A_{p}) \ge 3$). Both circumferences are measured on the same calibration scale, so the dimensionless aspect ratio is
\[
\Tstar = R/r = 5/7.
\]
\end{proof}

\begin{corollary}[Inverted-radius relation] \label{cor:R-less-r}
Under the cyclotomic-embedding calibration, $\Tstar < 1$, i.e.\ the major circle (additive flow) is shorter than the minor circle (multiplicative flow). In abstract $T^{2} = S^{1} \times S^{1}$ coordinates this is purely a labelling: the two $S^{1}$ factors are interchangeable up to which flow is identified with which circle. We chose the labels so that $S^{1}_{R}$ is the additive flow's circle and $S^{1}_{r}$ is the multiplicative flow's circle, which is the convention adopted in the Flatness Theorem (companion). Under that convention, the ratio is $5/7$ rather than $7/5$. (The standard $\mathbb{R}^{3}$ ring-torus embedding requires the major radius to exceed the minor; for the $\Zts$-torus a ring embedding therefore corresponds to relabelling, and the structural ratio is the unordered pair $\{5, 7\}$ with the additive/multiplicative identification fixing the order.)
\end{corollary}

\begin{remark}[On the ``narrow-major'' label]
Earlier drafts of this paper used the term ``narrow-major torus'' to describe the $\Tstar < 1$ regime. Since the standard $\mathbb{R}^{3}$ ring-torus convention takes $R > r$, this label is non-standard and we have replaced it with the explicit description above.
\end{remark}

\section{Companion appearances of $5/7$} \label{sec:companions}

The number $5/7$ arises in several places inside the authors' research program; we list them here for context. We are explicit about which are genuinely independent derivations and which are reformulations of the cyclotomic forcing of Theorem~\ref{thm:forced57}, since the convergence of multiple non-independent derivations on a target value is not, by itself, mathematical evidence.

\subsection{Reformulation 1: Cyclotomic reduction gap}

The same forcing of Theorem~\ref{thm:forced57} can be stated entirely at the level of cyclotomic fields: the gap between the smallest prime $p$ with $\Q(A_{p}) \subseteq \Q(\sqrt{D})$ for some $D$ (giving $p = 5$, $D = 5$) and the smallest prime $p$ with $\Q(A_{p})$ a proper degree-$3$ extension (giving $p = 7$) yields the ratio $5/7$. \emph{This is the same theorem in different language and is not an independent derivation.}

\subsection{Reformulation 2: Prime-$\pi$-$\varphi$ field bridge}

The compositum $\Q(\zeta_5, \zeta_7)$ has Galois group $(\Z/4\Z) \times (\Z/6\Z)$ over $\Q$; the totally real subfield $\Q(\zeta_5 + \zeta_5^{-1}, \zeta_7 + \zeta_7^{-1}) = \Q(\sqrt{5}, A_7)$ is degree $6$ over $\Q$, with the degree-$2$ subfield $\Q(\sqrt 5)$ at the prime $5$ and the degree-$3$ subfield $\Q(A_7)$ at the prime $7$. \emph{This is again the cyclotomic forcing of Theorem~\ref{thm:forced57}, stated as a Galois lattice; not independent.}

\subsection{Independent appearance 1: First-G law coprime window} \label{ssec:firstG}

The $\sinc^2$ resonance amplitude $\mathcal{R}(k,f) = \sin^2(\pi k f)/(k^2 \sin^2(\pi f))$ on $\Z/n\Z$ for $n = 10$ has a stability-window structure: at the smallest-prime-factor coprime window $W_p = (p-1)/p$, the value $W_5 = 4/5$ records the additive stability of the prime-$5$ component. The corresponding obstruction at $p = 7$ contributes the complementary fraction $1 - 1/7 = 6/7$. We stress that $4/5$ and $6/7$ do not directly combine arithmetically to $5/7$; rather, both arise from the same cyclotomic threshold (degree-$2$ at $5$, degree-$3$ at $7$) that fixes the ratio in Theorem~\ref{thm:forced57}. The point is that the threshold appears independently in the resonance-window analysis and in the torus-aspect analysis. See \cite{SandersGish-FirstG} for the resonance-window derivation; the connection to the present paper is via the common cyclotomic threshold, not via a direct arithmetic identity.

\subsection{Independent appearance 2: TSML/BHML harmony-cell ratio} \label{ssec:tsml}

The two composition tables on $\Zts$ have respective harmony-cell counts $|H(\mathrm{TSML})| = 73$ and $|H(\mathrm{BHML})| = 28$; the sum $73 + 28 = 101$ is prime. The ratio $73/101 = 0.7227\ldots$ is \emph{close to} but not equal to $5/7 = 0.7142\ldots$ (relative gap $\approx 1.2\%$). We do not claim $73/101 = 5/7$. The structural agreement at the $1\%$ level is a numerical observation that we record here as an open problem: whether the small discrepancy reflects a genuine asymptotic correction to the cyclotomic threshold under repeated table extension is not known. The earlier draft of this paper claimed exact agreement; that claim was incorrect and has been retracted.

\subsection{Status of the $5/7$ value across these appearances}

After the corrections above, the picture is:
\begin{itemize}
\item Theorem~\ref{thm:forced57} (this paper) is the single \emph{rigorous} cyclotomic forcing of $5/7$, conditional on the cyclotomic-embedding calibration of Definition~\ref{def:cyclo-cal}.
\item Reformulations~1--2 are equivalent statements of the same forcing.
\item Independent appearance~1 (\S\ref{ssec:firstG}) is a separate computation in which the same cyclotomic threshold appears, providing structural support but not an independent derivation of the ratio itself.
\item Independent appearance~2 (\S\ref{ssec:tsml}) is a numerical near-agreement at the $1\%$ level, recorded here as an open question rather than a derivation.
\end{itemize}
We have removed two earlier ``derivations'' (a BTQ operator balance and a $5/7 = (\sin(\pi/5)\sin(\pi/7))^{?}$ bridge) which a careful reader can verify do not constitute self-contained mathematical derivations of $5/7$.

\section{Generalization to other squarefree $n$} \label{sec:gen}

\begin{conjecture}[Generalized aspect ratio, restricted scope] \label{conj:gen}
Let $n \geq 6$ be squarefree, $n = p_1 \cdots p_k$ with $k \geq 2$, and assume that at least one prime divisor $p_i$ of $n$ satisfies $A_{p_i} \notin \Q$ and $\deg_{\Q}(A_{p_{i}}) \leq 2$. Define
\[
p_{\mathrm{closed}} := \min\{ p \mid n : A_p \notin \Q \text{ and } \deg_{\Q}(A_p) \leq 2 \},
\]
\[
p_{\mathrm{obstr}} := \min\{ p \text{ prime} : \deg_{\Q}(A_p) \geq 3 \} = 7.
\]
Then, under the cyclotomic-embedding calibration of Definition~\ref{def:cyclo-cal}, the torus aspect ratio is
\[
\Tstar(\Zn) = \frac{p_{\mathrm{closed}}}{p_{\mathrm{obstr}}} = \frac{p_{\mathrm{closed}}}{7}.
\]
\end{conjecture}

The hypothesis that some prime divisor $p_i \mid n$ has $A_{p_i}$ irrational and $\deg_{\Q}(A_{p_i}) \le 2$ is necessary: $A_2 = 0$ and $A_3 = 1$ are rational (degenerate), and at $p \geq 7$ the degree is already $\geq 3$ (no quadratic closure available among the prime divisors). The hypothesis is satisfied iff $5 \mid n$ (since $p = 5$ is the only prime with $A_p$ irrational of degree $\le 2$). The conjecture's domain is therefore exactly the squarefree multiples of $5$ that have at least one other prime factor.

\begin{proposition}[Necessity of the scope restriction] \label{prop:scope}
The function $n \mapsto p_{\mathrm{closed}}(n)$ is well-defined on $\{ n \in \mathbb{Z}_{\ge 1} : n \text{ squarefree}, \, 5 \mid n, \, n > 5 \}$ and undefined elsewhere among squarefree $n \geq 6$ with $k \geq 2$.
\end{proposition}

\begin{proof}
The set of primes $p$ with $A_p$ irrational and $\deg_{\Q}(A_p) \le 2$ is exactly $\{5\}$: for $p = 2, 3$, $A_p \in \Q$; for $p \ge 7$, $\deg_{\Q}(A_p) = (p-1)/2 \ge 3$. So $p_{\mathrm{closed}}(n)$ exists iff $5 \mid n$.
\end{proof}

\begin{itemize}
\item \emph{$n = 10$:} $p_{\mathrm{closed}} = 5$, $\Tstar = 5/7$ (Theorem~\ref{thm:forced57}).
\item \emph{$n = 15 = 3 \cdot 5$:} $p_{\mathrm{closed}} = 5$, conjectured $\Tstar = 5/7$.
\item \emph{$n = 35 = 5 \cdot 7$:} $p_{\mathrm{closed}} = 5$, conjectured $\Tstar = 5/7$ (here $p_{\mathrm{obstr}} = 7$ \emph{is} a prime divisor of $n$, but the conjecture remains $5/7$).
\item \emph{$n = 6 = 2 \cdot 3$:} no prime divisor with $\deg_{\Q}(A_{p}) \le 2$ irrational; conjecture is silent.
\item \emph{$n = 14 = 2 \cdot 7$:} no such prime divisor; conjecture is silent (this addresses the inconsistency in the earlier statement).
\item \emph{$n = 21 = 3 \cdot 7$:} no such prime divisor; conjecture is silent.
\end{itemize}

The conjecture is currently verified only for $n = 10$ (the present paper). The cases $n = 15$ and $n = 35$ are the most accessible candidates for extension; both require the four-structure embedding into a torus to be established by a generalization of the Flatness Theorem. A continuum-limit version (the asymptotic torus aspect ratio for $n \to \infty$ over $\{n : 5 \mid n, \, \text{squarefree}, \, n > 5\}$) would predict $\Tstar = 5/7$ universally on this restricted domain; this is also open.

\section{Open questions}

\begin{enumerate}[label=(\alph*)]
\item \emph{Conjecture~\ref{conj:gen}.} Verify the conjecture for $n = 15$ and $n = 35$, the next two cases in the conjecture's domain (squarefree multiples of $5$ with at least one additional prime factor).
\item \emph{Calibration-free derivation.} Theorem~\ref{thm:forced57} is conditional on the cyclotomic-embedding calibration imported from the Flatness Theorem. A self-contained derivation of the canonical calibration --- equivalently, a proof that the Flatness Theorem's torus admits a unique isometry class realizing the four ring structures --- would make $\Tstar = 5/7$ unconditional. This is the natural sequel.
\item \emph{Curvature of the ring torus.} Derive a closed-form expression for the Gaussian curvature of the $\Zn$-torus at each point $(\theta_A, \theta_M) \in T^2$ as a function of the additive and multiplicative phases. The standard ring-torus curvature has zero set at $\theta_M = \pi/2, 3\pi/2$; whether the cyclotomic data at $p = 7$ contributes an additional structural curvature term in the algebraic torus is open.
\item \emph{Connection to modular curves.} The sequence $\{T^2(\Zn)\}$ for squarefree $n$ in the conjecture's domain has cyclotomic structure tied to the cusps of the modular curve $X(N)$. The continuum limit (if it exists) is conjecturally a cusp neighborhood of a fundamental domain for $\mathrm{SL}(2,\Z)$ acting on the upper half-plane.
\item \emph{The $73/101$ near-agreement.} Determine whether the discrepancy $73/101 - 5/7 = 6/707$ between the TSML+BHML harmony-cell ratio and the cyclotomic threshold $5/7$ reflects an asymptotic correction (a $1/n$-type term in extended tables) or a genuine independence of the two numbers.
\end{enumerate}

\section*{Tier discipline}

We classify the central claims of this paper according to the standard PROVED / COMPUTED / STRUCTURAL RHYME / OPEN tier discipline used throughout the authors' research program:

\begin{itemize}
\item \textbf{PROVED.} Theorem~\ref{thm:forced57} (cyclotomic-calibrated $5/7$ aspect ratio on $\Z/10\Z$); Theorems~\ref{thm:R-forced} and \ref{thm:r-forced} (the major- and minor-radius selections); Lemma~\ref{lem:A7-irreducible} (irreducibility of $g(x) = x^{3} - x^{2} - 2x + 1$). The Galois group $A_{3} \cong \Z/3\Z$ and the discriminant $\disc(g) = 49$ are classical and verified.
\item \textbf{COMPUTED.} The minimal polynomials of $A_{2}, A_{3}, A_{5}, A_{7}$ over $\Q$; the discriminant $\disc(g) = 49$; the substitution bridge $h(x/2) = g(x)$ relating the $\cos(\pi/7)$ and $2\cos(\pi/7)$ polynomials; the 50-digit numerical zero $g(A_{7}) < 10^{-40}$. All reproduced at machine precision by the included script \texttt{verify\_J13.py} (6/6 PASS, runtime $< 5$ s).
\item \textbf{STRUCTURAL RHYME.} The numerical near-agreement $73 / 101 \approx 5/7$ (\S\ref{ssec:tsml}): an empirical $\approx 1.2\%$ proximity recorded as a structural observation, not a derivation. The First-G Law coprime windows $W_{5} = 4/5$ and $1 - 1/7 = 6/7$ (\S\ref{ssec:firstG}): independent appearances of the same cyclotomic threshold, providing structural support but not an independent derivation of $5/7$.
\item \textbf{OPEN.} (a) Conjecture~\ref{conj:gen} for $n \in \{15, 35\}$; (b) a calibration-free derivation (Remark~\ref{rem:calibration}); (c) curvature of the ring torus; (d) the modular-curve connection; (e) whether $73/101 - 5/7 = 6/707$ is an asymptotic correction or genuine independence.
\end{itemize}

\section*{Acknowledgments}

We thank M.~Gish for the additive flow / cyclotomic framing carried over from the Flatness Theorem, and C.~A.~Luther for discussions on the multiplicative flow's cubic obstruction at $p = 7$.

\begin{thebibliography}{9}

\bibitem{Lehmer}
D.~H.~Lehmer,
\emph{A note on trigonometric algebraic numbers},
Amer.\ Math.\ Monthly 40 (1933), 165--166.

\bibitem{WatkinsZeitlin}
W.~Watkins and J.~Zeitlin,
\emph{The minimal polynomial of $\cos(2 \pi / n)$},
Amer.\ Math.\ Monthly 100 (1993), no.\ 5, 471--474.

\bibitem{Niven}
I.~Niven,
\emph{Irrational Numbers},
Carus Mathematical Monographs 11, Mathematical Association of America, 1956.

\bibitem{DrapalWanless}
A.~Dr\'apal and I.~M.~Wanless,
\emph{Maximally nonassociative quasigroups},
J.\ Combin.\ Theory Ser.\ A \textbf{184} (2021), 105510.

\bibitem{Washington}
L.~C.~Washington,
\emph{Introduction to Cyclotomic Fields}, 2nd ed.,
Graduate Texts in Math.\ 83, Springer, 1997.

\bibitem{SandersGish-FirstG}
B.~R.~Sanders and M.~Gish,
\emph{First-G Law: Squarefree Stability of the Smallest-Prime-Factor Coprime Window},
submitted to \emph{Integers}, 2026.

\bibitem{SandersMayes-CL}
B.~R.~Sanders and B.~Mayes,
\emph{Crossing Lemma: Non-Associativity as Information Generation in Finite Magmas},
submitted to \emph{J.\ Combin.\ Theory Ser.\ A}, 2026.

\bibitem{SandersGish-Flatness}
B.~R.~Sanders and M.~Gish,
\emph{Flatness Theorem: The Forced $2 \times 2$ Torus on $\Z/10\Z$},
submitted to \emph{J.\ Pure Appl.\ Algebra}, 2026 (companion --- parent result).

\bibitem{SandersMayes-UOP}
B.~R.~Sanders and B.~Mayes,
\emph{The Universal Orthogonality Principle},
submitted to \emph{J.\ Number Theory}, 2026.

\end{thebibliography}

\end{document}
